\section*{Anexo E. Comparación con otros trabajos}
\phantomsection
\addcontentsline{toc}{subsection}{Anexo E. Comparación con otros trabajos}
\makeatletter\def\@currentlabel{E}\makeatother
\label{anx:trabajos}
\setcounter{table}{0}
\renewcommand{\thetable}{E.\arabic{table}}
\renewcommand{\theHtable}{E.\arabic{table}}
% fuente de todo el anexo: fichas de literatura/fichas_trabajos.json del 27.09.2026 (64 fichas, 54 obras), solo las de
% verificacion 'confirmado' o 'corregido'. Cada cifra se leyo en la tabla o el texto de la obra citada que
% indica su ficha; los porcentajes estan pasados a la convencion de la tesis (negativo = menos retraso,
% espera o tiempo de viaje). Las claves estan en referencias.bib desde el 27.09.2026.

Este anexo reúne los trabajos de control semafórico con los que la
Sección~\ref{sec:comparacion-literatura} compara los resultados del plan
de pruebas. La búsqueda se organizó en
cuatro frentes: bancos de prueba y revisiones, control de un cruce
aislado o de un corredor, redes con varios agentes, y trabajos de
América Latina junto con sistemas que operan en ciudades. De cada obra
se registraron la red, el simulador, la demanda, el algoritmo, las
entradas y salidas del agente, la referencia contra la que se mide, si
esa referencia estaba ajustada y las cifras con su métrica. Cada cifra
se releyó en la fuente, en el texto completo cuando estuvo disponible.
Cuando un trabajo reporta varios escenarios, la mediana se recalculó
con sus tablas. En once obras la relectura corrigió algún dato, como
los autores, la sede, la métrica o la referencia de una cifra. Una
investigación publicada como tesis y como artículos cuenta como una sola
obra: las cifras de MARLIN-ATSC salen de la tesis
\citep{ElTantawy2012_Tesis}, cuyos capítulos se publicaron también por
separado \citep{ElTantawy2012_ITSC,ElTantawy2014_Parametros,ElTantawy2013_MARLIN}.

Así se reunieron 54 obras. Quedaron fuera cuatro cuya cifra solo pudo
leerse en el resumen o en una fuente secundaria, entre ellas un
Q-learning para un cruce aislado \citep{Abdulhai2003_QLearning} y la
tesis de la UAQ revisada en el Capítulo~\ref{cap:estado-arte}
\citep{UAQ2024}, cuyo capítulo de resultados no estuvo accesible.
También quedó fuera un trabajo que mide el aprendizaje multiagente en un
simulador propio muy simplificado y contra un fijo de 5~s de verde
\citep{Mahato2025_MARLvsFijo}. Tampoco se usa el 73\,\% de mejora frente
a tiempo fijo que RESCO \citep{Ault2021_RESCO} atribuye a Mousavi et al.: en la fuente, esa
cifra es la reducción de la cola frente a otra red neuronal
\citep{Mousavi2017_PGDQN}. De las 49 obras restantes, siete no hacen una
comparación propia (dos simuladores o entornos y cinco revisiones) y 42
entran en las Tablas~\ref{tab:trabajos-comparados}
a~\ref{tab:trabajos-comparados-4}, una por frente y en orden
cronológico dentro de cada una.
% fuente: fichas con verificacion 'parcial' (fuera): Abdulhai2003_QLearning, Arel2010_MARL5 (sin clave
% en el .bib), CarrascoMorales2024_UAQ (clave UAQ2024) y Zamora2026_Peru (sin clave, sin texto completo).
% Mahato2025_MARLvsFijo: 'confirmado', pero con simulador en Pygame y fijo de 5 s (advertencia de su
% ficha). Mousavi: comparacion CMP4 de la ficha Ault2021_RESCO (corregida). Sin comparacion propia:
% Alegre2019_SUMORL, Zhang2019_CityFlow, Wei2019_Survey, Haydari2022_Survey, Noaeen2022_SLR,
% Yau2017_Survey y Greguric2020_Review. Obras con alguna ficha 'corregido': RESCO, LibSignal, Lioris,
% MARLIN-ATSC, Liang, Trinity, SURTRAC, Vivacity, SCOOT, FHWA (2008) y la ola verde de Lima.

Todas las cifras se expresan como cambio porcentual frente a la
referencia, con la convención de esta tesis: un valor negativo indica
menos retraso, espera o tiempo de viaje que la referencia, aunque el
trabajo original reporte la misma mejora como un porcentaje positivo.
Las cifras solo son comparables entre filas con la misma métrica. El
tiempo de viaje incluye el recorrido a velocidad libre, que ningún
control puede acortar, y por eso diluye el porcentaje. La espera
detenida es solo una parte del retraso y por eso infla el porcentaje. En los resultados de
X-Light, el mismo IPPO reduce el tiempo de viaje un 18.9\,\% y el
retraso un 40.4\,\% en la misma malla y frente al mismo fijo
\citep{Jiang2024_XLight}; en el piloto de
SURTRAC, el tiempo de viaje baja un 25.8\,\% y la espera un 40.6\,\%
\citep{Smith2013_Surtrac}. La métrica puede incluso invertir el signo:
en la misma malla y frente a Greedy, FMA2C reduce un 90.5\,\% el retraso
medido en el cruce y aumenta un 10.8\,\% el retraso del viaje
\citep{Ma2020_FMA2C}.
% fuente: X-Light CMP3 (IPPO en Grid4x4, 94.64 a 56.38 s de retraso; -18.9 % en tiempo de viaje);
% SURTRAC CMP0 y CMP1 (global, ponderado por volumen); FMA2C CMP2 y CMP3 (malla 4x4 con flujo creciente)

En las tablas de trabajos comparados, la columna Red da el número de
intersecciones semaforizadas de las redes evaluadas y si son reales o
sintéticas, y la columna Sim., el simulador; \emph{campo} indica una
medición en calles. En la referencia, un \emph{fijo sin tunear} usa
verdes iguales o arbitrarios, un \emph{fijo no declarado} es uno cuyo
ajuste el trabajo no describe y un \emph{fijo tuneado} se ajustó por
barrido o por programación lineal. El \emph{plan real} es el programa
que opera en la ciudad o el que trae el escenario, sin reoptimizar, y
\emph{otro RL} indica que la referencia es otro aprendiz. Varias líneas
con la misma referencia son escenarios distintos. Una cifra sin
aprendiz nombrado es la del mejor aprendiz que el trabajo reporta en ese
escenario, o la mediana de sus escenarios; con nombre, es la de ese
aprendiz. \emph{Máx.} marca una cifra que el trabajo da como máximo
(``hasta''), y \emph{sin cifra}, una comparación que el trabajo solo
muestra en figuras o describe sin números.

\begin{table}[H]
\setstretch{1}
\centering
\scriptsize
\setlength{\tabcolsep}{3pt}
\begin{tabular}{>{\raggedright\arraybackslash}p{1.9cm}>{\raggedright\arraybackslash}p{1.7cm}>{\raggedright\arraybackslash}p{1.25cm}>{\raggedright\arraybackslash}p{1.7cm}>{\raggedright\arraybackslash}p{2.8cm}>{\raggedright\arraybackslash}p{1.25cm}>{\raggedright\arraybackslash}p{1.8cm}}
\toprule
Trabajo & Red & Sim. & Algoritmo & Referencia & Métrica & Mejora (\%) \\
\midrule
Lioris et al. (2014) \citep{Lioris2014_MaxPressure} & 13, real & .Q & Max-Pressure \citep{Varaiya2013_MaxPressure}, sin RL
  & fijo tuneado (LP)
  & viaje
  & $-$23.6 \\
\addlinespace[2pt]
Genders y Razavi (2019) \citep{Genders2019_Framework} & 2, sintética & SUMO & DQN, DDPG
  & fijo tuneado\newline fijo tuneado\newline Webster adaptativo\newline Max-Pressure
  & viaje\newline viaje\newline viaje\newline viaje
  & $-$8.9\newline $-$1.3 DQN\newline $+$1.4\newline $+$22.0 \\
\addlinespace[2pt]
LIT (2019) \citep{Zheng2019_LIT} & 1 (12 cruces) y redes de 4 a 16, reales & SUMO & LIT (DQN)
  & fijo sin tunear\newline Webster\newline fijo sin tunear\newline Webster
  & viaje\newline viaje\newline viaje\newline viaje
  & $-$55.5\newline $-$34.0\newline $-$65.8 redes\newline $-$37.6 redes \\
\addlinespace[2pt]
RESCO (2021) \citep{Ault2021_RESCO} & 1 a 21, reales; mallas de 16, sintéticas & SUMO & IDQN, IPPO, MPLight, FMA2C
  & plan real\newline Max-Pressure\newline Greedy\newline otro RL (IDQN)
  & retraso\newline retraso\newline retraso\newline retraso
  & sin cifra\newline $-$11\newline $-$20\newline $-$2.1 IPPO \\
\addlinespace[2pt]
Cabrejas-Egea et al. (2021) \citep{CabrejasEgea2021_Comercial} & 1 (dos geometrías), sintética & Vissim & DDQN, DDDQN, A2C
  & fijo no declarado\newline MOVA (actuado)\newline MOVA (actuado)\newline MOVA (actuado)\newline SURTRAC
  & retraso\newline retraso\newline retraso\newline retraso\newline retraso
  & $-$85.2\newline $-$21.7\newline $-$48.0\newline $-$56.9\newline $-$26.6 \\
\addlinespace[2pt]
Efficient Pressure (2021) \citep{Wu2021_EfficientMP} & 12 y 16, reales & CityFlow & Efficient-CoLight (DQN)
  & fijo no declarado\newline Max-Pressure
  & viaje\newline viaje
  & $-$38.2\newline $-$3.7 \\
\addlinespace[2pt]
Advanced-XLight (2022) \citep{Zhang2022_AdvancedMP} & 12 a 192, reales & CityFlow & Advanced-CoLight (DQN)
  & fijo no declarado\newline Max-Pressure
  & viaje\newline viaje
  & $-$38.6\newline $-$8.6 \\
\addlinespace[2pt]
LibSignal (2024) \citep{Mei2024_LibSignal} & 1 a 16 en SUMO (196 en CityFlow), reales y sintéticas & SUMO y CityFlow & IDQN, IPPO, CoLight, MPLight y otros
  & fijo sin tunear\newline Max-Pressure\newline Max-Pressure
  & viaje\newline viaje\newline viaje
  & $-$55.1\newline $-$5.7\newline $+$18.5 IPPO \\
\addlinespace[2pt]
T-REX (2025) \citep{Nguyen2025_TREX} & 3 a 21, reales; malla de 16, sintética & SUMO & IDQN, IPPO, MPLight, FMA2C
  & plan real\newline plan real\newline fijo sin tunear\newline Max-Pressure
  & viaje\newline viaje\newline viaje\newline viaje
  & $-$17.7\newline $+$5.1 IPPO\newline $-$28.7 IDQN\newline $-$20.2 \\
\bottomrule
\end{tabular}
\caption[Trabajos comparados: bancos de prueba y control por presión]{Trabajos comparados, primer frente: bancos de prueba, comparaciones amplias y control por presión. Mejora: cambio porcentual frente a la referencia (negativo: menos retraso, espera o tiempo de viaje). Viaje: tiempo de viaje. Las convenciones de las columnas se explican en el texto.}
\label{tab:trabajos-comparados}
\end{table}
% fuente: fichas de cada obra en literatura/fichas_trabajos.json. LibSignal: mediana de 8 escenarios en SUMO (Tablas 4 y 7);
% Lioris: mediana de 7 pares O-D; Efficient Pressure: mediana de 5; Advanced-XLight: mediana de 6; T-REX:
% mediana de 3 redes reales (plan real) y de 4 escenarios (Max-Pressure), grilla 4x4 para el fijo sin
% tunear; el corredor de Ingolstadt queda fuera de las dos lineas de plan real (su tabla rotula mal las
% filas de los otros controles; IPPO si se lee ahi: -10.9, ficha Nguyen2025_TREX CMP2 de la 2a ficha); RESCO: validacion en las mallas 4x4 (Fig. 4), IPPO frente a IDQN es la mediana de 6 tareas.

\begin{table}[H]
\setstretch{1}
\centering
\scriptsize
\setlength{\tabcolsep}{3pt}
\begin{tabular}{>{\raggedright\arraybackslash}p{1.9cm}>{\raggedright\arraybackslash}p{1.7cm}>{\raggedright\arraybackslash}p{1.25cm}>{\raggedright\arraybackslash}p{1.7cm}>{\raggedright\arraybackslash}p{2.8cm}>{\raggedright\arraybackslash}p{1.25cm}>{\raggedright\arraybackslash}p{1.8cm}}
\toprule
Trabajo & Red & Sim. & Algoritmo & Referencia & Métrica & Mejora (\%) \\
\midrule
MARLIN-ATSC (2012) \citep{ElTantawy2012_Tesis} & 1, 5 y 59, reales & Paramics & Q-learning, independiente y coordinado
  & Webster, 1 cruce\newline actuado, 1 cruce\newline Webster, 5 cruces\newline actuado, 5 cruces\newline plan real, 59 cruces
  & retraso\newline retraso\newline retraso\newline retraso\newline retraso
  & $-$35.5\newline $-$25.1\newline $-$21.4\newline $-$18.7\newline $-$37.6 \\
\addlinespace[2pt]
Genders y Razavi (2016) \citep{Genders2016_DQTSC} & 1, sintética & SUMO & DQN (CNN)
  & otro RL
  & retraso
  & $-$82.4 \\
\addlinespace[2pt]
Gao et al. (2017) \citep{Gao2017_DQN} & 1, sintética & SUMO & DQN
  & fijo no declarado\newline LQF
  & retraso\newline retraso
  & $-$86 máx.\newline $-$47 máx. \\
\addlinespace[2pt]
Mousavi et al. (2017) \citep{Mousavi2017_PGDQN} & 1, sintética & SUMO & gradiente de política y DQN (CNN)
  & fijo sin tunear\newline otro RL
  & retraso\newline retraso
  & sin cifra\newline $-$72.8 \\
\addlinespace[2pt]
IntelliLight (2018) \citep{Wei2018_IntelliLight} & 1 (sintética y 24 reales, una a una) & SUMO & DQN
  & fijo tuneado\newline SOTL
  & viaje\newline viaje
  & $-$40.3\newline $-$22.1 \\
\addlinespace[2pt]
FRAP (2019) \citep{Zheng2019_FRAP} & 1 (18 cruces con demanda real) & SUMO & FRAP (DQN)
  & fijo tuneado\newline Webster
  & viaje\newline viaje
  & $-$64.1\newline $-$58.1 \\
\addlinespace[2pt]
Liang et al. (2019) \citep{Liang2019_3DQN} & 1, sintética & SUMO & 3DQN
  & fijo sin tunear
  & espera
  & $-$25.7 \\
\addlinespace[2pt]
PressLight (2019) \citep{Wei2019_PressLight} & 3 a 16, reales; 6 a 20, sintéticas & CityFlow & DQN con presión
  & fijo sin tunear\newline onda verde\newline Max-Pressure
  & viaje\newline viaje\newline viaje
  & $-$52.1\newline $-$50.6\newline $-$42.3 \\
\addlinespace[2pt]
Wang et al. (2019) \citep{Wang2019_EventData3DQN} & 1, sintética & SUMO & 3DQN
  & Webster\newline actuado
  & retraso\newline retraso
  & $-$21.2\newline $-$10.1 \\
\addlinespace[2pt]
Ault et al. (2020) \citep{Ault2020_Interpretable} & 1, real & SUMO & DQN regulable, PPO
  & actuado\newline actuado, demanda alta
  & retraso\newline retraso
  & $-$19.4 máx.\newline peor, PPO \\
\bottomrule
\end{tabular}
\caption[Trabajos comparados: cruce aislado y corredor]{Trabajos comparados, segundo frente: cruce aislado y corredor. Mismas convenciones que la Tabla~\ref{tab:trabajos-comparados}.}
\label{tab:trabajos-comparados-2}
\end{table}
% fuente: MARLIN-ATSC, Tablas 5-8 (perfil uniforme), 6-12 (demanda actual) y 7-5 de la tesis; IntelliLight,
% datos reales de Jinan; FRAP, mediana de 18 cruces; PressLight, mediana de 10 escenarios.

\begin{table}[H]
\setstretch{1}
\centering
\scriptsize
\setlength{\tabcolsep}{3pt}
\begin{tabular}{>{\raggedright\arraybackslash}p{1.9cm}>{\raggedright\arraybackslash}p{1.7cm}>{\raggedright\arraybackslash}p{1.25cm}>{\raggedright\arraybackslash}p{1.7cm}>{\raggedright\arraybackslash}p{2.8cm}>{\raggedright\arraybackslash}p{1.25cm}>{\raggedright\arraybackslash}p{1.8cm}}
\toprule
Trabajo & Red & Sim. & Algoritmo & Referencia & Métrica & Mejora (\%) \\
\midrule
S2R2L (2019) \citep{Guo2019_S2R2L} & 16, sintética & SUMO & DQN descentralizado
  & Max-Pressure\newline Max-Pressure con verde máximo ajustado
  & retraso\newline retraso
  & $-$34.6\newline $-$19.9 \\
\addlinespace[2pt]
CoLight (2019) \citep{Wei2019_CoLight} & 12 a 196, reales; 36, sintética & CityFlow & DQN con atención
  & fijo sin tunear\newline Max-Pressure
  & viaje\newline viaje
  & $-$25.2\newline $-$12.7 \\
\addlinespace[2pt]
MA2C (2020) \citep{Chu2020_MA2C} & 30, real; 25, sintética & SUMO & MA2C
  & Greedy\newline otro RL (IA2C)\newline Greedy
  & retraso\newline retraso\newline retraso en el cruce
  & $-$20.8\newline $-$30.9\newline $-$64.1 \\
\addlinespace[2pt]
MPLight (2020) \citep{Chen2020_MPLight} & 2510, real; 16, sintética & CityFlow & DQN con presión
  & fijo sin tunear\newline Max-Pressure
  & viaje\newline viaje
  & $-$47.5\newline $-$14.4 \\
\addlinespace[2pt]
FMA2C (2020) \citep{Ma2020_FMA2C} & 30, real; 16, sintética & SUMO & MA2C jerárquico
  & Greedy\newline Greedy\newline otro RL (MA2C)\newline Greedy
  & retraso\newline retraso\newline retraso\newline retraso en el cruce
  & $-$6.3 Mónaco\newline $+$10.8 malla\newline $-$55.7\newline $-$90.5 malla \\
\addlinespace[2pt]
AttendLight (2020) \citep{Oroojlooy2020_AttendLight} & 1 (112 cruces), reales y sintéticas & CityFlow & atención y gradiente de política
  & fijo sin tunear\newline Max-Pressure
  & viaje\newline viaje
  & $-$46\newline $-$39 \\
\addlinespace[2pt]
EPyMARL en RESCO (2022) \citep{Ault2022_EPyMARL} & 3 a 21, reales & SUMO & nueve MARL (IPPO, MAPPO, QMIX y otros)
  & plan real, corredor de 3 cruces
  & retraso
  & ninguno lo supera \\
\addlinespace[2pt]
X-Light (2024) \citep{Jiang2024_XLight} & 8 a 25, reales y sintéticas & SUMO & PPO con \textit{transformer}
  & fijo sin tunear\newline fijo sin tunear
  & viaje\newline viaje
  & $-$28.8\newline $-$27.0 IPPO \\
\addlinespace[2pt]
CoSLight (2024) \citep{Ruan2024_CoSLight} & 8 a 29, reales y sintéticas & SUMO & selección de colaboradores
  & fijo sin tunear\newline Max-Pressure\newline Max-Pressure\newline fijo sin tunear, red de Nanshan
  & viaje\newline viaje\newline viaje\newline viaje
  & $-$27.3\newline $-$13.8\newline $-$5.7\newline $+$2.0 IPPO, $+$2.1 MAPPO \\
\bottomrule
\end{tabular}
\caption[Trabajos comparados: redes con varios agentes]{Trabajos comparados, tercer frente: redes con varios agentes. Retraso en el cruce: retraso medido solo en la intersección, no en todo el viaje. Mismas convenciones que la Tabla~\ref{tab:trabajos-comparados}.}
\label{tab:trabajos-comparados-3}
\end{table}
% fuente: CoLight frente a Max-Pressure, mediana de los 5 escenarios de la ficha (-6.6, -12.7, -10.7,
% -29.6, -19.4), calculada para este anexo; MPLight, mediana de 5; X-Light y CoSLight frente al fijo,
% mediana de 5 escenarios; IPPO y MAPPO de CoSLight, red de Nanshan; FMA2C y MA2C, Monaco salvo la malla.

\begin{table}[H]
\setstretch{1}
\centering
\scriptsize
\setlength{\tabcolsep}{3pt}
\begin{tabular}{>{\raggedright\arraybackslash}p{1.9cm}>{\raggedright\arraybackslash}p{1.7cm}>{\raggedright\arraybackslash}p{1.25cm}>{\raggedright\arraybackslash}p{1.7cm}>{\raggedright\arraybackslash}p{2.8cm}>{\raggedright\arraybackslash}p{1.25cm}>{\raggedright\arraybackslash}p{1.8cm}}
\toprule
Trabajo & Red & Sim. & Algoritmo & Referencia & Métrica & Mejora (\%) \\
\midrule
SCOOT (1999) \citep{DETR1999_SCOOT} & varias ciudades, reales & campo & SCOOT, sin RL
  & fijo actualizado\newline fijo típico\newline actuado aislado
  & retraso\newline retraso\newline retraso
  & $-$12\newline $-$20\newline $-$23 \\
\addlinespace[2pt]
FHWA (2008) \citep{Koonce2008_TSTM} & síntesis de EE.~UU. & -- & tiempos actualizados, sin RL
  & operación actual (proyección)
  & retraso
  & $-$15 a $-$40 \\
\addlinespace[2pt]
SURTRAC (2013) \citep{Smith2013_Surtrac} & 9, real & campo & planificación en línea, sin RL
  & actuado coordinado\newline actuado coordinado
  & viaje\newline espera
  & $-$25.8\newline $-$40.6 \\
\addlinespace[2pt]
Higuera (2016) \citep{Higuera2016_Uniandes} & 6, real (Bogotá) & SUMO & Q-learning coordinado
  & plan real
  & espera
  & $-$62.5 \\
\addlinespace[2pt]
ASCT, FHWA (2017) \citep{FHWA2017_EDC1ASCT} & síntesis de EE.~UU. & -- & control adaptativo, sin RL
  & temporización existente
  & viaje
  & $\leq -10$ \\
\addlinespace[2pt]
City Brain (2018) \citep{Alibaba2018_CityBrain} & un distrito, real & campo & IA sin detallar
  & situación previa
  & velocidad
  & $+$15 \\
\addlinespace[2pt]
Ola verde de Lima (2019) \citep{Sovero2019_OlaVerde} & siete avenidas, real (Lima) & campo & onda verde por franja horaria, sin RL
  & situación previa (cifra municipal)
  & viaje
  & $-$10 \\
\addlinespace[2pt]
Borges et al. (2021) \citep{Borges2021_Brasilia} & 2, real & Vissim & \textit{deep Q-learning}
  & plan real
  & retraso
  & $-$28.4 \\
\addlinespace[2pt]
Montoya (2021) \citep{Montoya2021_PUCP} & 10, real (Lima) & SUMO & algoritmo genético, sin RL
  & plan real\newline onda verde manual
  & espera\newline espera
  & $-$84.7\newline $-$78.6 \\
\addlinespace[2pt]
Vivacity (2021) \citep{Velasco2021_Vivacity} & 5, real & campo & RL sin detallar
  & incumbente\newline actuado (simulación)
  & viaje\newline espera
  & $-$30\newline $-$45 \\
\addlinespace[2pt]
Trinity (2022) \citep{Chakrabarti2022} & 1 y 2, sintéticas & SUMO & DQN, A2C, PPO
  & fijo sin tunear
  & espera
  & sin cifra \\
\addlinespace[2pt]
Uniandes (2023) \citep{Chavez2023} & 4, sintética & SUMO & Q-learning, SARSA, DQN, PPO
  & fijo sin tunear\newline fijo sin tunear
  & espera\newline espera
  & $-$63.5\newline $-$35.8 Q-learning \\
\addlinespace[2pt]
Tarquino (2023) \citep{Tarquino2023_Lima} & 7 (5 con semáforo), real (Lima) & Aimsun & actuado con detectores, sin RL
  & plan real\newline plan real
  & retraso\newline viaje
  & $-$35.4\newline $-$21.2 \\
\addlinespace[2pt]
Green Light (2023) \citep{Google2023_GreenLight} & más de 100 ciudades, reales & campo & IA que recomienda planes
  & plan real
  & paradas
  & $-$30 máx. \\
\bottomrule
\end{tabular}
\caption[Trabajos comparados: América Latina y sistemas en operación]{Trabajos comparados, cuarto frente: América Latina y sistemas en operación. Campo: medición en calles. En velocidad, un valor positivo es mejora. Mismas convenciones que la Tabla~\ref{tab:trabajos-comparados}.}
\label{tab:trabajos-comparados-4}
\end{table}
% fuente: SCOOT, TAL 7/99 (estimaciones promedio y Worcester); FHWA (2008), proyeccion de NTOC (rango
% de 15 a 40 %); SURTRAC, Tabla 1 (global); Higuera, Tabla 8-1 (Q-BR); Montoya, Tabla 4.2 (nivel alto);
% Tarquino, Tabla 6 (escenario 1); Uniandes, DQN y Q-learning (resumen y Sec. 8).

Las Tablas~\ref{tab:trabajos-entradas} a~\ref{tab:trabajos-entradas-3}
describen, para las obras con aprendizaje por refuerzo que publican su
formulación, qué observa el agente, qué decide, con qué recompensa
aprende y con qué métricas se evalúa. No incluyen los controles sin
aprendizaje ni el informe de Vivacity \citep{Velasco2021_Vivacity}, que
no describe su formulación. La penúltima fila resume una revisión
sistemática de 160 artículos sobre redes de dos o más cruces
\citep{Noaeen2022_SLR}, y la última describe esta tesis.

\begin{table}[H]
\setstretch{1}
\centering
\scriptsize
\setlength{\tabcolsep}{3pt}
\begin{tabular}{>{\raggedright\arraybackslash}p{2.0cm}>{\raggedright\arraybackslash}p{3.8cm}>{\raggedright\arraybackslash}p{2.4cm}>{\raggedright\arraybackslash}p{2.4cm}>{\raggedright\arraybackslash}p{1.8cm}}
\toprule
Trabajo & Estado (entradas) & Acción & Recompensa & Métricas de salida \\
\midrule
SUMO-RL (2019) \citep{Alegre2019_SUMORL}
  & fase, indicador de verde mínimo, densidad y cola por carril
  & siguiente fase verde, con ámbar
  & cambio de la espera acumulada
  & no reporta \\
\addlinespace[2pt]
Genders y Razavi (2019) \citep{Genders2019_Framework}
  & densidad y cola por carril, fase
  & siguiente fase (DQN) o duración del verde (DDPG)
  & $-$retraso de los vehículos del cruce
  & viaje, cola, retraso \\
\addlinespace[2pt]
LIT (2019) \citep{Zheng2019_LIT}
  & vehículos por carril, fase
  & mantener o cambiar
  & $-$suma de colas
  & viaje \\
\addlinespace[2pt]
RESCO (2021) \citep{Ault2021_RESCO}
  & por carril en 200~m: cola, vehículos que se acercan, espera, velocidad, llegadas y salidas (cada algoritmo usa una parte)
  & combinación de fases compatibles cada 10~s
  & $-$espera (IDQN, IPPO), $-$presión (MPLight), de vecindario (FMA2C)
  & retraso, viaje, espera, cola \\
\addlinespace[2pt]
Cabrejas-Egea et al. (2021) \citep{CabrejasEgea2021_Comercial}
  & cola por carril (m), etapa
  & etapa a activar
  & $-$suma de colas
  & retraso, cola \\
\addlinespace[2pt]
Efficient Pressure (2021) \citep{Wu2021_EfficientMP}
  & fase, presión eficiente
  & siguiente fase, de 15~s
  & $-$presión o $-$cola
  & viaje \\
\addlinespace[2pt]
Advanced-XLight (2022) \citep{Zhang2022_AdvancedMP}
  & fase, presión eficiente y vehículos en marcha en el rango efectivo
  & siguiente fase, mínimo 15~s
  & $-$presión o $-$cola
  & viaje \\
\addlinespace[2pt]
LibSignal (2024) \citep{Mei2024_LibSignal}
  & vehículos por carril y fase; presión y fase (MPLight)
  & siguiente fase, acíclica
  & vehículos en espera, espera (IPPO) o presión
  & viaje, cola, retraso, \textit{throughput} \\
\addlinespace[2pt]
T-REX (2025) \citep{Nguyen2025_TREX}
  & por carril: fase, vehículos que se acercan, detenidos, espera y velocidad; presión (MPLight)
  & fase compatible, de 10~s
  & reducción del retraso acumulado; $-$presión (MPLight)
  & viaje, retraso, espera, cola, robustez \\
\addlinespace[2pt]
MARLIN-ATSC (2012) \citep{ElTantawy2012_Tesis}
  & colas y llegadas por fase; acciones de los vecinos en el modo coordinado
  & siguiente fase, en secuencia variable
  & reducción del retraso acumulado
  & retraso, \textit{throughput}, paradas, cola, viaje, CO$_2$ \\
\addlinespace[2pt]
Genders y Razavi (2016) \citep{Genders2016_DQTSC}
  & presencia y velocidad por celda (imagen), fase
  & una de 4 fases
  & cambio del retraso acumulado
  & \textit{throughput}, cola, viaje, retraso \\
\addlinespace[2pt]
Gao et al. (2017) \citep{Gao2017_DQN}
  & posición y velocidad por celdas de 8~m, fase
  & verde O--E o N--S durante 10~s
  & reducción del tiempo de permanencia
  & retraso \\
\bottomrule
\end{tabular}
\caption[Formulación de los trabajos con aprendizaje por refuerzo]{Formulación de los trabajos con aprendizaje por refuerzo: qué observa el agente (estado), qué decide (acción), con qué señal aprende (recompensa) y con qué métricas se evalúa. Mismo orden que las Tablas~\ref{tab:trabajos-comparados} a~\ref{tab:trabajos-comparados-4}, con el entorno SUMO-RL al inicio.}
\label{tab:trabajos-entradas}
\end{table}

\begin{table}[H]
\setstretch{1}
\centering
\scriptsize
\setlength{\tabcolsep}{3pt}
\begin{tabular}{>{\raggedright\arraybackslash}p{2.0cm}>{\raggedright\arraybackslash}p{3.8cm}>{\raggedright\arraybackslash}p{2.4cm}>{\raggedright\arraybackslash}p{2.4cm}>{\raggedright\arraybackslash}p{1.8cm}}
\toprule
Trabajo & Estado (entradas) & Acción & Recompensa & Métricas de salida \\
\midrule
Mousavi et al. (2017) \citep{Mousavi2017_PGDQN}
  & imagen del simulador (128 $\times$ 128, 4 cuadros)
  & verde N--S o E--O
  & cambio del retraso acumulado
  & retraso, cola \\
\addlinespace[2pt]
IntelliLight (2018) \citep{Wei2018_IntelliLight}
  & por carril: cola, vehículos y espera; imagen de posiciones; fase actual y siguiente
  & mantener o cambiar
  & suma ponderada de cola, retraso, espera, cambios, vehículos que pasan y viaje
  & viaje, cola, índice de retraso \\
\addlinespace[2pt]
FRAP (2019) \citep{Zheng2019_FRAP}
  & vehículos por movimiento, fase
  & fase del siguiente intervalo (8)
  & $-$cola media por movimiento
  & viaje \\
\addlinespace[2pt]
Liang et al. (2019) \citep{Liang2019_3DQN}
  & posición y velocidad en 60 $\times$ 60 celdas
  & $\pm$5~s a una fase del ciclo siguiente
  & cambio de la espera acumulada entre ciclos
  & espera \\
\addlinespace[2pt]
PressLight (2019) \citep{Wei2019_PressLight}
  & fase, vehículos en los carriles de salida y por tramo en los de entrada
  & una de 4 fases, acíclica
  & $-$presión
  & viaje \\
\addlinespace[2pt]
Wang et al. (2019) \citep{Wang2019_EventData3DQN}
  & ocupación de detectores de lazo en ventanas de 60~s, estado de las señales
  & una de 4 fases, acíclica
  & vehículos que entran menos espera ponderada
  & retraso, \textit{throughput}, cola, velocidad, paradas \\
\addlinespace[2pt]
Ault et al. (2020) \citep{Ault2020_Interpretable}
  & señal actual; por carril: vehículos que se acercan, espera, detenidos y velocidad
  & par de fases compatibles, mínimo 3~s
  & reducción del retraso acumulado
  & retraso \\
\addlinespace[2pt]
S2R2L (2019) \citep{Guo2019_S2R2L}
  & ocupación y cola por carril, etapa, fracción de verde transcurrida; las de los vecinos
  & etapa cada 5~s, verde de 10 a 60~s
  & cambio de la espera acumulada, ponderada con la de los vecinos
  & retraso, cola, combustible \\
\addlinespace[2pt]
CoLight (2019) \citep{Wei2019_CoLight}
  & fase, vehículos por carril; atención sobre los vecinos
  & fase del siguiente intervalo
  & $-$suma de colas
  & viaje \\
\addlinespace[2pt]
MA2C (2020) \citep{Chu2020_MA2C}
  & vehículos a menos de 50~m y espera del primero, por carril; estados y políticas de los vecinos
  & fase cada 5~s
  & $-$(cola + espera), con la de los vecinos descontada
  & cola, retraso en el cruce y en el viaje, velocidad \\
\addlinespace[2pt]
MPLight (2020) \citep{Chen2020_MPLight}
  & fase, presión de 12 movimientos
  & una de 8 fases cada 10~s
  & $-$presión
  & viaje, \textit{throughput} \\
\addlinespace[2pt]
FMA2C (2020) \citep{Ma2020_FMA2C}
  & vehículos que llegan y espera por carril; flujo por dirección en el nivel regional
  & una de 5 fases; meta de flujo regional
  & como MA2C, más una recompensa regional
  & cola, retraso en el cruce y en el viaje, velocidad \\
\bottomrule
\end{tabular}
\caption[Formulación de los trabajos con aprendizaje por refuerzo (continuación)]{Continuación de la Tabla~\ref{tab:trabajos-entradas}.}
\label{tab:trabajos-entradas-2}
\end{table}

\begin{table}[H]
\setstretch{1}
\centering
\scriptsize
\setlength{\tabcolsep}{3pt}
\begin{tabular}{>{\raggedright\arraybackslash}p{2.0cm}>{\raggedright\arraybackslash}p{3.8cm}>{\raggedright\arraybackslash}p{2.4cm}>{\raggedright\arraybackslash}p{2.4cm}>{\raggedright\arraybackslash}p{1.8cm}}
\toprule
Trabajo & Estado (entradas) & Acción & Recompensa & Métricas de salida \\
\midrule
AttendLight (2020) \citep{Oroojlooy2020_AttendLight}
  & por carril: vehículos en marcha en 3 tramos y en espera, agregados con atención
  & fase, verde mínimo 10~s
  & $-$presión
  & viaje \\
\addlinespace[2pt]
EPyMARL en RESCO (2022) \citep{Ault2022_EPyMARL}
  & el de RESCO
  & fases compatibles
  & $-$espera en 200~m, común a todos
  & retraso \\
\addlinespace[2pt]
X-Light (2024) \citep{Jiang2024_XLight}
  & cola, fase, ocupación, flujo y detenidos
  & una de 8 fases cada 15~s
  & suma ponderada de cola, espera, retraso y presión
  & viaje, retraso \\
\addlinespace[2pt]
CoSLight (2024) \citep{Ruan2024_CoSLight}
  & por movimiento: vehículos, cola, fase y flujo; datos de los colaboradores
  & fase cada 15~s
  & común a todos los cruces
  & viaje, retraso \\
\addlinespace[2pt]
Higuera (2016) \citep{Higuera2016_Uniandes}
  & hora del día, fase, cola máxima y espera por acceso; las del vecino (Q-BR)
  & una de 2 fases, verde mínimo 20 o 14~s
  & $-$suma de cuadrados de cola y espera
  & cola, espera, velocidad, viaje, CO$_2$ \\
\addlinespace[2pt]
Borges et al. (2021) \citep{Borges2021_Brasilia}
  & flujo, velocidad y tiempo de cada sensor
  & verde o rojo por cruce cada 10~s, mínimo 30~s
  & flujo de cada cruce y entre los dos
  & retraso, velocidad, viaje, cola \\
\addlinespace[2pt]
Trinity (2022) \citep{Chakrabarti2022}
  & fase; por carril, espera y densidad de detenidos
  & fase en cada paso
  & $-$espera
  & espera \\
\addlinespace[2pt]
Uniandes (2023) \citep{Chavez2023}
  & fase, verde mínimo cumplido, densidad y cola por carril
  & una de 4 fases verdes
  & cambio del retraso acumulado
  & ``eficiencia'' frente al fijo \\
\addlinespace[2pt]
Revisión de 160 artículos \citep{Noaeen2022_SLR}
  & cola 38\,\%, fase 11\,\%, vehículos 10\,\%, posición 6\,\%, velocidad 6\,\%
  & por fase 62\,\%, por ciclo 32\,\%
  & cola 30\,\%, retraso 13\,\%, espera 9\,\%, vehículos 6\,\%
  & retraso 20\,\%; viaje, espera y cola 12\,\% cada una \\
\midrule
Esta tesis
  & tabulares: detenidos por aproximación en 4 niveles y fase verde; DQN, IPPO y MAPPO: vector de 31 componentes (53 en la malla de P5 y P6) con cola detenida y vehículos en movimiento por aproximación, fase, tiempo en verde, verde mínimo cumplido, identificador y mensajes de los vecinos (Tabla~\ref{tab:obs-u4})
  & mantener o pasar a la siguiente fase verde, cada 5~s; verde de 10 a 60~s
  & $-$(retraso $+$ 0.01 $\times$ espera), sumado por segundo desde la decisión anterior y dividido entre 5 (Ecuación~\ref{eq:recompensa-u4}); término de \textit{spillback} en IPPO spill; colas en la ablación
  & retraso por vehículo (\texttt{timeLoss}, primaria); espera detenida, colas, paradas, \textit{throughput}, CO$_2$ y NO$_x$; en el análisis de robustez, retraso total con la espera para entrar a la red (Sección~\ref{sec:robustez}) \\
\bottomrule
\end{tabular}
\caption[Formulación de los trabajos con aprendizaje por refuerzo (continuación)]{Continuación de la Tabla~\ref{tab:trabajos-entradas}. La fila de la revisión da la frecuencia de cada elemento en los estudios que revisa, en porcentaje; la última fila describe esta tesis.}
\label{tab:trabajos-entradas-3}
\end{table}
% fuente: campos estado, accion, recompensa y metricas de cada ficha; fila de la revision: Noaeen2022_SLR
% (preprint de engrXiv); fila de esta tesis: Capitulo III (sec:formulacion, sec:u4, Tabla tab:obs-u4 y sec:u3);
% metrica primaria del resultado principal: timeloss_prom (plan2/ESPEC.md, veredicto del 25.09); retraso total solo
% en la robustez (plan2/robustez_veredicto.csv).

La mejora reportada depende sobre todo de la referencia. Con una cifra
por obra y por tipo de referencia, solo de trabajos con aprendizaje por
refuerzo, sin informes de fabricantes y solo en retraso, espera o
tiempo de viaje, la mediana es de $-46.0$\,\% frente a un fijo sin
tunear o no declarado (15 obras), de $-35.7$\,\% frente a un fijo
tuneado, un plan de Webster o el plan real (10 obras), de $-20.7$\,\%
frente al actuado (4 obras) y de $-15.5$\,\% frente a Max-Pressure u
otra heurística adaptativa (17 obras). Buena parte de la distancia
entre la primera categoría y la última la cubre una heurística sin
entrenamiento: frente al mismo fijo sin tunear o no declarado, Max-Pressure ya reduce
el tiempo de viaje un 35.8\,\% en un trabajo \citep{Wu2021_EfficientMP}
y un 48.5\,\% en otro \citep{Mei2024_LibSignal}. Las cifras más bajas
aparecen donde la referencia está mejor ajustada. Frente a un fijo de
verde uniforme tuneado por barrido y evaluado con 32 semillas, el mejor
aprendiz reduce el tiempo de viaje un 8.9\,\%
\citep{Genders2019_Framework}. Frente al plan real, IPPO aumenta el
tiempo de viaje en el corredor de Colonia y en la región de Ingolstadt
\citep{Nguyen2025_TREX}, y ningún algoritmo de EPyMARL supera al
plan en el corredor de tres cruces \citep{Ault2022_EPyMARL}.
% fuente de las medianas: calculadas para este anexo sobre las fichas (literatura/gen.py y
% literatura/build.py, 27.09.2026; se rehacen a mano con esta lista). Reglas: solo obras con RL y sin informes de
% fabricante (fuera Vivacity, Green Light y City Brain); una cifra por obra y por tipo de referencia, que es
% la mediana de sus lineas de ese tipo en las tablas (cada linea ya es la mediana de escenarios que da la
% ficha, si la da). No entran las lineas que muestran un aprendiz distinto del mejor (IPPO en LibSignal,
% T-REX, X-Light y CoSLight; DQN en Genders y Razavi 2019; Q-learning en Uniandes; PPO en Ault et al.
% 2020), las de otro RL ni las de retraso en el cruce; la linea de T-REX con IDQN en la malla si entra, porque
% es la unica de ese tipo en su ficha. Frente a una heuristica, Max-Pressure si la obra lo compara en ese
% escenario y si no Greedy, LQF, SOTL o SURTRAC (en RESCO, Max-Pressure en la malla Grid y Greedy en la
% malla Avenues); si la obra ajusta la heuristica, la version ajustada.
% Fijo sin tunear o no declarado (n = 15, mediana -46.0, rango -86 a -25.2): LibSignal -55.1; T-REX -28.7
%   (grilla, IDQN); LIT -60.65 (mediana de -55.5 y -65.8); Cabrejas-Egea -85.2; Efficient Pressure -38.2;
%   Advanced-XLight -38.6; Gao -86 (max.); Liang -25.7; PressLight -52.1; CoLight -25.2; MPLight -47.5;
%   AttendLight -46; X-Light -28.8; CoSLight -27.3; Uniandes -63.49. Sin cifra: Mousavi, Trinity.
% Fijo tuneado, Webster o plan real (n = 10, mediana -35.65, rango -62.49 a -8.9): Genders y Razavi 2019
%   -8.9; LIT -35.8 (mediana de -34.0 y -37.6); T-REX -17.7; MARLIN-ATSC -35.5 (mediana de -35.5, -21.4 y
%   -37.6); IntelliLight -40.3; FRAP -61.1 (mediana de -64.1 y -58.1); PressLight -50.6 (onda verde);
%   Wang -21.2; Higuera -62.49; Borges -28.4. Sin cifra: RESCO, EPyMARL.
% Actuado (n = 4, mediana -20.65, rango -48.0 a -10.1): Cabrejas-Egea -48.0 (mediana de -21.7, -48.0 y
%   -56.9); MARLIN-ATSC -21.9 (mediana de -25.1 y -18.7); Wang -10.1; Ault et al. 2020 -19.4 (max.).
% Max-Pressure u otra heuristica (n = 17, mediana -15.5, rango -47 a +22.0): RESCO -15.5 (mediana de -11 y
%   -20); Genders y Razavi 2019 +22.0; Cabrejas-Egea -26.6 (SURTRAC); Efficient Pressure -3.7; LibSignal
%   -5.7 (SUMO); Advanced-XLight -8.6; T-REX -20.2; Gao -47 (LQF, max.); IntelliLight -22.1 (SOTL);
%   PressLight -42.3; S2R2L -19.9 (ajustado); CoLight -12.7; MA2C -20.8 (Greedy); MPLight -14.4; FMA2C
%   +2.25 (mediana de -6.3 y +10.8); AttendLight -39; CoSLight -9.75 (mediana de -13.8 y -5.7).
% Max-Pressure sin RL frente al fijo sin tunear: Efficient Pressure CMP3 (-35.8, mediana de 5) y LibSignal
% CMP5 (-48.5, mediana de 8 en SUMO). IPPO frente al plan real en T-REX: +50.2 (corredor de Colonia),
% -3.3 (region de Colonia) y +5.1 (region de Ingolstadt) (ficha Nguyen2025_TREX, CMP2); en el corredor de
% Ingolstadt, -10.9 (2a ficha Nguyen2025_TREX, CMP2). EPyMARL: Cologne Corridor, sin cifras en el trabajo.
